\section{The Open Account}

The argument can now be laid out whole. The debt is real: existence received and presently sustained is the first and chief cause of indebtedness. Settlement is impossible: nothing the creature could tender is not already the creditor's, and no rendering reaches equality. The impossibility is the debt's species, not its defect: what is owed is not a sum but an acknowledgment---equality of wills where equality of things cannot be. The virtue that renders it is annexed to justice: sober as debt, humble as sign, engaging the whole embodied creature, which knows and speaks through its body. The rendering profits only the renderer, who is the one party capable of profit; and it must be free, because a consent extracted is not a consent. Whatever else remains obscure about religion, its skeleton is this.

One temptation remains, and it is the subtlest, because it wears the argument's own clothes. A creature persuaded of the whole argument can still quietly retain the ledger picture at one remove: very well, the debt cannot be \emph{closed}---but perhaps it can be \emph{serviced}; perhaps a life of punctual acknowledgment, vows kept and knees bent, accumulates a kind of standing, until the worshipper, without ever saying so, comes to regard himself as paid up---square with heaven at last, self-possessed in the moral order as no creature is in the ontological one. This is recoil's final disguise: not the refusal of worship but the \emph{use} of worship to end the relation worship acknowledges. And it fails as every version fails: the acts were never payments; each was itself received---the borrowed breath spent on naming the lender is still borrowed breath, and the naming, too, is on loan. Aquinas's serene remark closes the door: the debt of gratitude is the kind that grows by being paid, ``and it is not unreasonable if the obligation of gratitude has no limit.''\endnote{\emph{Summa theologiae} II-II, q.~106, a.~6, ad~2, checked 2026-07-24 as in section~2 above. The present paragraph's target---fulfilled religious acts construed as installments toward a discharged, self-possessed standing before God---is this essay's own formulation, offered as project synthesis and attributed to no source; the underlying principle (no aggregate of acts closes the creature's relation to its sustaining source) follows from ST I, q.~104, a.~1 as argued in Part I.} Every act of the due return, performed, enlarges what there is to acknowledge---one more gift taken with the eyes open---and so the account stands further open, not nearer closed. A worshipper who understands this stops experiencing the open account as arrears. The account is open the way a window is open, not the way a wound is: it is the relation itself, seen from the side of action, and to want it closed is to want the relation over---which was recoil's desire from the beginning.

The psalm knew. Its most quoted line about servitude comes wrapped around an image of release, and the two are not fighting: ``O Lord, for I am thy servant: I am thy servant, and the son of thy handmaid. Thou hast broken my bonds: I will sacrifice to thee the sacrifice of praise.''\endnote{Psalm 115:16--17, Douay--Rheims, checked 2026-07-24 as in section~1 above. The verses' juxtaposition of avowed servanthood and broken bonds is the textual ground for this paragraph's reading; the resolution offered---the bond to the unfailing principle as the breaking of every other bond---is this essay's synthesis, in the sense of \latin{religare} from II-II, q.~81, a.~1 and of CCC 2097 (``The worship of the one God sets man free from turning in on himself'').} I am thy servant; thou hast broken my bonds: in any human institution the sentences contradict. Here they explain each other. The bond to the unfailing principle---\latin{religare}, the etymology that binds---is the breaking of every other bond, because a creature that has acknowledged the one total claim can no longer be totally claimed by anything else: not by the idols that accept payment in full, not by the State that may not touch the debt, not by the self's own project of getting square. The psalmist says it in two verses that end, like this series' every trail, at public thanksgiving.

And justice's last act, here as everywhere, is to point past itself. The analysis of gratitude left an arrow embedded in the argument: the limitless debt ``flows from charity''---from love, which is why paying enlarges it. Aquinas's map of the virtues says the same structurally---the theological virtues \emph{command} religion---and the Catechism's first-commandment paragraph opened with the same order: it is charity that ``leads us to render to God what we as creatures owe him in all justice.'' Follow the arrow and the entire apparatus of debt begins to look penultimate. Justice can establish that the return is owed and describe its acts; it cannot say what the relation is \emph{for}---why the fountain wanted drinkers, why the singer wanted a song that could answer. A benefactor, Aquinas noticed, is not finally after repayment; the exceeding, gratis-upon-gratis structure of gratitude already behaves less like accounting than like friendship. What the giver wants, the whole tradition ventures, is not the payment but the payer. That claim---that the relation of total dependence is aimed at something scandalously mutual; that the creature is not only sustained and obligated but \emph{desired}; that the end of the whole architecture is not acknowledgment but union---is not a claim justice can underwrite or the ascent can reach. It is the matter of charity, and of revelation, and of Part III.

Part I ended with the discovery that the falling was being carried. This essay has added the second, which is the first made practical: the carried creature can answer. Not repay---answer: take the cup knowingly, say the giver's name aloud, kneel until the body has taught the mind, hand over under signs a self that was never self-possessed, keep the vow in public, and do it all freely or not at all. The due return is not the price of the gift; nothing was ever the price of the gift. It is the gift's completion in the only place completion could occur---in the receiver, who alone can change, and whose change is this: the gift at last received \emph{as} a gift, by someone who knows it, says so, and stays.
